\documentclass[language=english]{../../Template/TemplateLecture}

\renewcommand{\LectureName}{\textsc{Functional Analysis and PDE}}
\renewcommand{\ProfName}{\textsc{Prof. Dr. Stefan Müller}}
\renewcommand{\Semester}{Winter 2024/25}
\renewcommand{\mName}{Daniel Emse}

\begin{document}
% \setcounter{chapter}{0}

%%%%%%%%%%%%%%%%%%%%%%%%%%%%%%%%%%%%%%%%%%%%%%%%%%%%%%%%%%%%%%%%%%%%%%%%%%%%%%%%%%%%%%%%%%%%%%%%%%%%
%%%%%%%%%%%%%%%%%%%%%%%%%%%%%%%%%%%%%%%%%% 2024 / 04 / 09 %%%%%%%%%%%%%%%%%%%%%%%%%%%%%%%%%%%%%%%%%%
%%%%%%%%%%%%%%%%%%%%%%%%%%%%%%%%%%%%%%%%%%%%%%%%%%%%%%%%%%%%%%%%%%%%%%%%%%%%%%%%%%%%%%%%%%%%%%%%%%%%
\chapter{Overview/Teaser}

\newLecture{09.10.2024}

\begin{enumerate}
    \item What is FA?
    \begin{itemize}
        \item Analysis on infinite-dimensional spaces

        \emph{examples:} $C^0(X), C^\infty(U \subseteq \IR^n), \ell_p = \set{a: \IN \to \IR \mid \sum_{n = 0}^\infty \abs{a_n}^p < \infty} ~(1 \leq p \leq \infty)$
        \item Combination of geometric, analytic and (linear) algebraic ideas
    \end{itemize}
    \item What is FA good for?
    \begin{itemize}
        \item Geometry/Topology with manifolds, bundles, loop space
        \item PDEs, scalar products
        \item Algorithms/Numerics (approxmiation theory; errors)
        \item infinite-dimensional representations
        \item Mathematical Physics / Quantum Mechanics; spectral theory
        \item Probability Theory (infinite products)
    \end{itemize}
    \item Why combine PA \& PDE?
    \begin{itemize}
        \item $u: U \subseteq \IR^n \to \IR$ as a $C^2$ function:
        \begin{equation}
            \Delta u(x) = \sum_{i = 1}^n \frac{\partial^2 u}{\partial x_i^2}.
        \end{equation}
        with $-\Delta u = f$.

        \emph{Classical approach}: Given a potential $f \in C^2$, does there exist a solution and is that solution unique?

        \emph{Modern approach}: $\Delta: X \to Y$ linear map between function spaces. e.g. $X = C^2(U), Y = C^0(U)$. Existance of solutions for all $f \in C^0$ is equiv to surjectivity of $f\Delta$ and uniqueness is equivalent to injectivity.

        Fun fact: $C^2$ is a bad space for $\Delta$. But $\Delta$ is a good operator on many function spaces, e.g. $C^{2,\alpha}$ with $0 < \alpha < 1$, so $D^2 u$ is $\alpha$-Hölder continuous, or for $W^{2,p}$ as Sobolev spaces (where $D^2 u \in L^p$ for $1 < p < \infty$), where $D^2$ is the weak derivative
    \end{itemize}
    \item \enquote{Big things}
    \begin{itemize}
        \item Spaces (of founctions) and convergence (norm, scalarproduct)
        \item (Bounded) Linear operators, Hahn-Banach theorem, uniform boundedness, \ldots
        \item Hilbert spaces, projections
        \item Spectral theory (LA: $A: \IR^n \to \IR^n$ symmetric, so $\Span{Ax, y} = \Span{x, Ay} \Rightarrow \exists$ orthogonal basis of eigenvectors; we want something similar in infinite dimensions which makes diagonazing easier, and we can easily compute the inverse or powers of the operator)
    \end{itemize}
\end{enumerate}

\chapter{Structures}
Generally, all vector spaces will be with $\IK = \IR$ or $\IK = \IC$.
\begin{itemize}
    \item on \enquote{infinite-dimensional} spaces
    \item Convergence and neighbourhoods, topological space
    \item Metric space (metric, triangle inequality, gives a \emph{distance})
    \item Normed space (vector space, $\norm{\lambda x} = \abs{\lambda} \cdot \norm{x}$ for $\lambda \in \IR$, gives a \emph{length})
    \item Scalar product space ($\norm{x} = \sqrt{\Span{x,x}}$, gives a \emph{length and an angle}) 
\end{itemize}

\section{Topological spaces}

Let $X$ a set and $2^X$ the power set.

\begin{defi}{Topological space, Topology}{topSpace}
    A \emph{topological space}\index{Space!Topological -}\index{Topological Space} $(X, \cT)$ is a pair of a set $X$ and subset $\cT$ of $2^X$ with the following properties:
    \begin{enumerate}[$(i)$]
        \item $\emptyset \in \cT, X \in \cT$
        \item finite intersections: $U, V \in \cT \Rightarrow U \cap V \in \cT$
        \item arbitrary unions: $U_\alpha \in \cT ~\forall \alpha \in A \Rightarrow \bigcup_{\alpha \in A} U_\alpha \in \cT$.
    \end{enumerate}
    The set $\cT$ is called a \emph{Topology on $X$}\index{Topology} and the elements in $\cT$ are often called \emph{open sets}\index{open set}\index{set!open -}. A set $C$ is called \emph{closed}\index{closed set}\index{set!closed -} if $X \setminus C$ is open.

    If $\cT_1$ and $\cT$ are topologies on $X$, then $\cT_1$ is called \emph{finer} than $\cT_2$ or $\cT_2$ is called \emph{coarser} than $\cT_1$ if $\cT_1 \supset \cT_2$.
\end{defi}

\begin{bsp}
    \begin{enumerate}[$(i)$]
        \item The finest topology is $\cT = 2^X$.
        \item The coarsest topology is $\cT = \set{\emptyset, X}$.
        \item Standard topology on $\IR$: $U \in \cT_{st} \iff \forall x \in U \exists \varepsilon > 0: \intervaloo{x-\varepsilon, x+\varepsilon} \subseteq U$. We will usually use this topology.
        \item Let $X = \IR$ and $\cT = \set{U \subseteq \IR \mid U = \emptyset \text{ or $\IR \setminus U$ is countable}}$. Then, the closed sets are countable\footnote{We call a set \emph{countable}\index{countable}\index{set!countable} if it is empty, finite or has an bijection to $\IN$.} or $\IR$. 
        \item Let $(X, \cT)$ be a topological space. Then, for $A \subseteq X$, the \emph{relative topology}\index{topology!relative -} $\cT_A$ on $A$ is defined by $\cT_A = \set{U \cap A \mid U \in \cT}$.
        \item Let $B$ be a set and $\cT_\beta$ a topology on $X$ for all $\beta \in B$. Then, $\cT := \bigcap_{\beta \in B} \beta$ defines a topology on $X$.
        
        Given any $\cS \subseteq 2^X$,
        \begin{equation}
            \bigcap_{\substack{\cT~\text{top.}\\\cT \supseteq \cS}} \cT
        \end{equation}
        is a topology and is called the \emph{coarsest topology containing $\cS$}.
        \item \emph{Product topology} of $(X, \cT_X), (Y, \cT_Y)$ as topological spaces is definied with 
        \begin{equation}
            \cS := \set{U \times Y \mid U \in \cT_X} \cup \set{X \times V \mid V \in \cT_Y}.
        \end{equation}
        The product topology is $\cT_\cS$, the coarsest topology containing $\cS$.
        
        The product topology is the coarsest topology so that the projections $\pi_1(x,y) = x$ and $\pi_2(x,y) = y$ are \enquote{continuous}.
    \end{enumerate}
\end{bsp}

\begin{defi}{Interior, Closure, Boundary}{interorClosureBoundary}
    Let $(X, \cT)$ be a topological space and $A \subseteq X$. We define the \emph{interior}\index{set!interior} $A^\circ$, the \emph{closure}\index{set!closure} $\overline{A}$ and the \emph{boundary}\index{set!boundary} $\partial A$ by 
    \begin{equation}
        A^\circ = \bigcup_{\substack{U \subseteq A \\ U \in \cT}} U, \quad \overline{A} = \bigcap_{\substack{C \supseteq A \\ X \setminus C \in \cT}} C, \quad \partial A = \overline{A} \setminus A^\circ.
    \end{equation}
\end{defi}

\begin{defi}{Dense subsets, separable topological space}{denseSubsets}
    $A \subseteq X$ is \emph{dense}\index{set!dense} if $\overline{A} = X$. The space $X$ is \emph{separable}\index{separable}\index{Topological Space!separable -} if there exists a countable, dense
    subset.
\end{defi}

\begin{defi}{Open neighbourhood}{openNeighbourhood}
    Let $x \in X$. Then, $U \subseteq X$ is an \emph{open neighbourhood of $x$}\index{set!open neighbourhood}\index{open neighbourhood} if $U \in \cT$ and $x \in U$.  
\end{defi}

\begin{defi}{Continuous maps, Homeomorphisms}{continuosity}
    Let $(X, \cT_X), (Y, \cT_Y)$ topological spaces. Then, $f: X \to Y$ is called \emph{continuous}\index{continuous} if
    \begin{equation}
        \forall V \supseteq \set{x} \in \cT_Y: \quad f^{-1}(V) = \set{x \in X \mid f(x) \in V} \in \cT_X.
    \end{equation}
    We call $f$ a \emph{homeomorphism}\index{homeomorphism} if $f$ is bijective and both $f$ and $f^{-1}$ are continuous.
\end{defi}

If $\cT = 2^X$, and $f: X \to Y$ is continuous.

If $\cT = \set{\emptyset, \cT}$, then every constant map is continuous.




\newLecture{11.10.2024}

\begin{defi}{Convergence}{convergence}
    Let $(X, \cT)$ be a topological space and $x: \IN \to X$ a sequence. Let $x^* \in X$. We say that \emph{$x$ converges to $x^*$}\index{convergence} if for \textbf{every} neighbourhood $U$ of $x$, the set $\set{j \in \IN \mid x_j \notin U}$ is finite.
\end{defi}

\begin{beispiel}
    \begin{enumerate}[$(i)$]
        \item If $\cT = \set{\emptyset, X}$, every sequence converges to every point $x^* \in X$.
        \item If $\cT = 2^X$, take $U = \set{x^*}$, so only sequences which are constant with up to finitely many exceptions converge.
    \end{enumerate}
\end{beispiel}

\begin{defi}{Hausdorff space}{hausdorffSpace}
    A topological space is called a \emph{Hausdorff space}\index{Hausdorff}\index{Hausdorff!- space}\index{Space!Hausdorff -} if 
    \begin{enumerate}[$(i)$]
        \item sets consisting of a single point are closed,
        \item for different points $x \neq y$, there exist disjoint open neighbourhoods: $\exists U_x, U_y \in \cT$ with $x \in U_x, y \in U_y$ and $U_x \cap U_y = \emptyset$.
    \end{enumerate}
\end{defi}

\begin{prop}
    If $(X, \cT)$ is Hausdorff, then every sequence has at most one limit point. \hfill $\qed$
\end{prop}

\begin{defi}{Compactness}{compactness}
    A subset $K$ of a topological space $(X, \cT)$ is \emph{compact}\index{compact}\index{set!compact -} if every open cover of $K$ contains a \textbf{finite} subcover: If $A$ is an index set and $U_\alpha \in \cT$ for every $\alpha \in A$ and 
    \begin{equation}
        \bigcup_{\alpha \in A} U_\alpha \supseteq K,
    \end{equation}
    then there exists a \textbf{finite} set $A' \subseteq A$ so that 
    \begin{equation}
        \bigcup_{\alpha \in A'} U_\alpha \supseteq K.
    \end{equation}
\end{defi}

\begin{lem}{Continuous maps compact to compact}{compactUnderContinuous}
    If $f: (X, \cT_X) \to (Y, \cT_Y)$ is continuous and $K \subseteq X$ is compact in $X$, then $f(X)$ is compact in $Y$.
\end{lem}

\begin{proof}
    Follows directly from the definitions.
\end{proof}

\begin{thm}{Continuous maps with compact domain attain max and min}{compactToMaxMin}
    If $f: (X, \cT_X) \to \IR$ is continuous and $K \subseteq X$ is compact, then $f$ attains its minimum and maximum on $K$, i.e. $\exists \underline{x}, \overline{x} \in K$ so that $f(\underline{x}) = \inf_{x \in K} f(x), f(\overline{x}) = \sup_{x \in K} f(x)$.
\end{thm}

\begin{proof}
    $f(K) \subseteq \IR$ is compact, the rest is easy.
\end{proof}

\begin{defi}{Connectedness}{connectedness}
    Let $(X, \cT)$ be a topological space. We call $X$ \emph{connected}\index{Topological Space!connected -}\index{connected}\index{Space!connected -} if it cannot be written as a disjoint union of two open subsets. A subset $A \subseteq X$ is connected if $(A, \cT_A)$ is connected.
\end{defi}

\begin{aufgabe}
    Show that intervals in $\IR$ are connected.
\end{aufgabe}

\begin{thm}{Intermediate value theorem}{intermediateValue}
    If $f: (X, \cT_X) \to \IR$ is continuous and $(X, \cT_X)$ is connected, and there exist $a, b \in X$ so that $f(a) < 0, f(b) > 0$. Then, there exists an $x \in X$ with $f(x) = 0$.
\end{thm}

\begin{proof}
    Assume $f(x) \neq 0$ for all $x \in X$. We know that $U := f^{-1}(\intervaloo{0,\infty})$ is open, and $X \setminus U = f^{-1}(\intervaloc{-\infty,0}) = f^{-1}(\intervaloo{-\infty,0})$ is also open. Then, $X = U \cup (X \setminus U) \Rightarrow X$ not connected. $\lightning$
\end{proof}

\section{Metric spaces}
\subsection{Definition and examples}

Let $X$ be a set and $d: X \times X \to \intervalco{0,\infty}$.

\begin{defi}{Metric space, metric}{metricSpace}
    $(X,d)$ is called a \emph{metric space}\index{Space!metric -}\index{metric space} and $d$ is called \emph{metric}\index{metric} if
    \begin{enumerate}[$(i)$]
        \item $d(x,y) = 0 \Leftrightarrow x = y$,
        \item $\forall x,y: d(x,y) = d(y,x)$,
        \item (triangle inequality)\index{triangle inequality} $d(x,y) \leq d(x,z) + d(z,y)$ for all $x,y,z \in X$.\footnote{\enquote{Where the triangle inequality is not true is on airline flights. Sometimes, it is cheaper to fly from Köln/Bonn to Mallorca first before flying to mainland Spain later.} -- Stefan Müller, 11.10.2024}
    \end{enumerate}
    If only $(ii)$ and $(iii)$ are fulfilled, we call $d$ a \emph{semimetric} or \emph{pseudometric}\index{semimetric}\index{pseudometric}\index{metric!semi-}\index{metric!pseudo-}, and with $x \sim y :\Leftrightarrow d(x,y) = 0$, we get a metric $\overline{d}$ on the equivalence classes.
\end{defi}

\begin{lem}{Composition of concave map and metric is metric}{compositionConcaveMetric}
    Assume $\psi: \intervalco{0,\infty} \to \intervalco{0,\infty}$ is continuous, non-decreasing, \textbf{concave} and $\psi(0) = 0, \psi(x) > 0$ for $x > 0$. If $d$ is a metric on $X$, so is $\psi \circ d$.
\end{lem}

\begin{proof}
    It is enough and easy to check that $\psi(a+b) \leq \psi(a) + \psi(b)$. 
\end{proof}

\begin{beispiel}
    Define $\psi(t) := \frac{t}{1+t} = 1 - \frac{1}{1+t}$. Then, $\psi \in C^0(\intervalco{0,\infty}) \cap C^2(\intervaloo{0,\infty})$. Non-decreasing and concave are equivalent to $\psi' \geq 0$ and $\psi'' \leq 0$.
\end{beispiel}

\begin{beispiel}
    \begin{enumerate}[$(i)$]
        \item Let $X = \IR^n$ and $d_\infty(x,y) = \max_{1 \leq i \leq n} \abs{x_i - y_i}$.
        \item $X = \IR^n, d_p(x,y) := \left(\sum_{i=1}^n \abs{x_i - y_i}^p\right)^{\frac{1}{p}}$ for $1 \leq p < \infty$,
        \item $\ell_p = \set{a: \IN \to \IR \mid \left(\sum_{i=1}^\infty \abs{x_i}^p\right)^{\frac{1}{p} < \infty}}$ and $d_p(a,b) := \left(\sum_{i=1}^\infty \abs{a_i - b_i}^p\right)^{\frac{1}{p}}$ for $1 \leq p < \infty$
        \item $\ell_\infty := \set{a: \IN \to \IR \mid \sup_{i \in \IN} \abs{a_i} < \infty}$ with $d_\infty(a,b) := \sup_{i \in \IN} \abs{a_i - b_i}$
        \item $X = \IR$ with $d(x,y) = \frac{\abs{x-y}}{1+\abs{x-y}} = \psi(\abs{x-y})$ for $\psi(t) = \frac{t}{1+t}$. By Lemma~\ref{lem:compositionConcaveMetric}, $d$ is a metric on $\IR$. Fun fact: $0 \leq d(x,y) < 1$.
        \item Metric on the space of all sequences:
        \begin{equation}
            d(a,b) = \sum_{i = 0}^\infty 2^{-i} \frac{\abs{a_i - b_i}}{1 + \abs{a_i - b_i}} \leq \sum_{i = 0}^\infty 2^{-i} \leq 2.
        \end{equation}
        \item Let $X$ be a set and $(Y, d_Y)$ a metric space and $f: X \to Y$ \textbf{injective}. Then, we can define the \emph{pullback metric}\index{metric!pullback -}\index{pullback metric} $d_X(x,y) := d_Y(f(x),f(y))$. $d_X$ is a metric on $X$.
        \item $\overline{\IR} := \IR \cup \set{-\infty,\infty}$. Then, 
        \begin{equation}
            f: \overline{\IR} \to \intervalcc{-1,1}, x \mapsto \begin{cases}
                -1 & x = -\infty \\ \frac{x}{1+\abs{x}} & x \in \IR \\ 1 & x = \infty
            \end{cases}
        \end{equation}
        is bijective. Then, $d_{\overline{\IR}}(x,y) = \abs{f(x)-f(y)}$ is a pullback metric.
        \item Shortest path on a graph
    \end{enumerate}
\end{beispiel}

\begin{bemerkung}
    Let $(X,d)$ be a metric space, let $A, B \subseteq X$, let $x \in X$ and let $r > 0$. Further distances:
    \begin{description}
        \item[Diameter of a subset] 
        \begin{equation}
            \sup\set{d(x,y) \mid x,y \in A}
        \end{equation} 
        (with the convention $\diam(A) = \infty$ if $\set{d(x,y) \mid x,y \in A}$ is unbounded)
        \item[Distance between two sets]
        \begin{equation}
            \dist(A,B) = \inf\set{d(x,y) \mid x \in A, y \in B}
        \end{equation}
        \item[Distance between a point and a set]
        \begin{equation}
            \dist(x,A) = \inf\set{d(x,y) \mid y \in B} = \dist(\set{x},A)
        \end{equation}
        \item[$r$-neighbourhood of a set]
        \begin{equation}
            B_r(A) := \set{x \in X \mid \dist(x,A) < r}
        \end{equation}  
        \item[$r$-ball of a point]
        \begin{equation}
            B_r(x) = B(x,r) = \set{y \in X \mid d(y,x) < r}
        \end{equation}
        \item[Hausdorff distance of two sets]\index{Hausdorff distance}\index{set!Hausdorff distance between -s}
        \begin{equation}
            d_H(A,B) := \inf\set{r > 0 \mid A \subseteq B_r(B) \text{ and } B \subseteq B_r(A)}
        \end{equation}
        and $d_H(\set{x},\set{y}) = d(x,y)$.   
    \end{description}
\end{bemerkung}

\begin{bemerkung}
    It is easy to show that
    \begin{equation}
        d_H(A,B) = \max\left(\sup_{x \in A} \dist(x,B), \sup_{y \in B} \dist(y,A)\right).
    \end{equation}
\end{bemerkung}

\begin{proof}
    Exercise.
\end{proof}

\begin{prop}[Topology induced by a metric]
    Let $(X,d)$ be a metric space. Let $\cT_d$ consist of all sets $U$ with the following property:
    \begin{equation}
        \forall x \in U \exists \varepsilon > 0: B_\varepsilon(x) \subseteq U.
    \end{equation}
    Then, $\cT_d$ is a topology on $X$.
\end{prop}

\begin{proof}
    Clear. For the intersection, take the ball with the smaller radius.
\end{proof}

\begin{lem}{Hausdorff distance is metric on closed and bounded sets}{hausdorffDistanceMetric}
    Let $(X,d)$ be a Hausdorff space. Then, $d_H$ defines a metric on
    \begin{equation}
        \cC := \set{A \subseteq X \mid \emptyset \neq A \text{ is closed and bounded}}.
    \end{equation}
\end{lem}

\begin{proof}
    Exercise.
\end{proof}

\newLecture{16.10.2024}

\subsection{Convergence and continuity in metric spaces}

\begin{prop}
    Let $(X,d)$ be a metric space and $x : \IN \to X$ a sequence, and $x^* \in X$. Then, the following two statements are equivalent:
    \begin{enumerate}[$(i)$]
        \item $x$ converges to $x^*$ in $(X,\cT_d)$,
        \item $\lim_{n \to \infty} d(x_n, x^*) = 0$.
    \end{enumerate}
\end{prop}

\textbf{Moral.} \enquote{In metric spaces, anything can be done by sequences}

\begin{prop}
    Let $(X,d)$ be a metric space and let $A \subseteq X$. Then, the following two statements are equivalent:
    \begin{enumerate}[$(i)$]
        \item $A$ is closed, so $X \setminus A \in \cT_d$,
        \item $A$ is \emph{sequentially closed}\index{set!sequentially closed -}\index{set!closed -!sequentially -}, so for every sequence $x: \IN \to A$ which converges to $x^* \in X$, one has $x^* \in A$.
    \end{enumerate}
    Moreover, the closure of any set $A$ is given by the \emph{sequential closure}\index{sequential closure}, this means 
    \begin{equation}
        \overline{A} = \set{x^* \in X \mid \exists x: \IN \to A, \text{$x$ converges to $x^*$}}.
    \end{equation}
\end{prop}

\begin{proof}
    Easy exercise.
\end{proof}

\begin{prop}
    Let $(X, d_X), (Y, d_Y)$ be two metric spaces and $f: X \to Y$. Then, the following three statements are equivalent:
    \begin{enumerate}[$(i)$]
        \item $f$ is continuous in as a map $f: (X, \cT_{d_X}) \to (Y, \cT_{d_Y})$,
        \item For every $x\in X, \varepsilon > 0$ there exists a $\delta > 0$ so that $d_X(z,x) < \delta \Rightarrow d_Y(f(z), f(x)) < \varepsilon$,
        \item For every sequence $x$ convergent to $x^*$, also $f(x)$ converges to $f(x^*)$.
    \end{enumerate}
\end{prop}

We know that every metric space is Hausdorff. \emph{Question:} If $(X, \cT)$ is a Hausdorff space, does there exist a metric $d: X \times X \to \intervalco{0,\infty}$ such that $\cT = \cT_d$? Or shorter, is $X$ metrizable?

If $X$ is second countable (so there exist countable neighbourhoods), $X$ is metrizable, and then, $X$ ist first countable.

\begin{defi}{Equivalent metrices, stronger metrices}{equivalentMetrices}
    Let $d_1$ and $d_2$ be metrices on $X$. 
    
    We say \emph{$d_1$ is stronger than $d_2$} if $\cT_{d_1} \supset \cT_{d_2}$.\index{metric!stronger -}

    We say \emph{$d_1$ is equivalent to $d_2$} if $\cT_{d_1} = \cT_{d_2}$.\index{metric!equivalent -}
\end{defi}

\begin{prop}
    The following statements are equivalent:
    \begin{enumerate}[$(i)$]
        \item $d_1$ is stronger than $d_2$,
        \item The identity map $x \mapsto x$ for $(X,d_1) \to (X,d_2)$ is continuous,
        \item Every sequence that converges in $d_1$ also converges in $d_2$,
        \item $\forall x \in X \forall \varepsilon > 0 \exists \delta > 0$ so that $d_1(x,y) < \delta \Rightarrow d_2(x,y) < \varepsilon$.
    \end{enumerate}
\end{prop}

\begin{beispiel}\leavevmode
    \begin{enumerate}[$(i)$]
        \item Let $X = \IR$ and $d_1(x,y) = \abs{x - y}, d_2 = \abs{\arctan(x) - \arctan(y)}$. Then, $d_1$ and $d_2$ are equivalent (Exercise 1.4).
        \item Let $\ell_1 := \set{x: \IN \to \IR \mid \sum_{k = 0}^\infty \abs{x_k} < \infty}$ with $d_1(x,y) := \sum_{k = 0}^\infty \abs{x_k - y_k}$ and $d_\infty(x,y) = \sup_{k \in \IN} \abs{x_k - y_k} \leq d_1(x,y)$. We see that $d_1$ is stronger than $d_\infty$, but not equivalent: See the map 
        \begin{equation}
            \IN \to (\IN \to \IR), \quad j \mapsto x^{(j)}_k := \begin{cases}
                0 & k < j-1 \text{ or } k > 2j \\
                \frac{1}{k} & j \leq k \leq 2j.
            \end{cases}
        \end{equation} 
        Then, $\abs{x_k^{(j)}} \leq \frac{1}{j}$ with $d_\infty(x^{(j)},0) \leq \frac{1}{j} \to 0$ as $j \to \infty$, but $d_1(x_k^{(j)},0) = \sum_{k = 0}^\infty x_k^{(j)} = \sum_{k = j}^{2j} \frac{1}{k} \geq j\frac{1}{2j} = \frac{1}{2}$.
    \end{enumerate}
\end{beispiel}

\begin{defi}{Separable metric space}{separable}
    A metric space $(X,d)$ is called \emph{separable}\index{metric space!separable -}\index{separable} if $(X, \cT_d)$ is separable.
\end{defi}


\subsection{Completeness}

\begin{defi}{Cauchy sequence, completeness}{cauchySeqCompleteness}\index{Cauchy sequence}\index{sequence!Cauchy -}\index{complete}\index{metric space!complete -}
    Let $(X,d)$ be a metric space.
    \begin{enumerate}[$(i)$]
        \item A sequence $x: \IN \to X$ is a \emph{Cauchy sequence} if $\forall \varepsilon > 0 \exists k_0$ so that for all $k, \ell \geq k_0: d(x_k, x_\ell) < \varepsilon$,
        \item $(X,d)$ is called \emph{complete} if every Cauchy sequence converges.
    \end{enumerate}
\end{defi}

\begin{beispiel}\leavevmode
    \begin{enumerate}[$(i)$]
        \item $\IQ$ equipped with the standard metric is not complete.
        \item $(\ell_p, d_p)$ are complete for all $1 \leq p \leq \infty$.
        \item $\IR$ equipped with the standard metric is complete, but not with the metric $\tilde{d}(x,y) := \abs{\arctan(x) - \arctan(y)}$, although they induce the same topology (Exercise 1.4).
        \item $\ell_1 = \set{x: \IN \to \IR \mid \sum_{k = 0}^\infty \abs{x_k} < \infty}$ with $d_1$ is complete (see $(ii)$), but $(\ell_1, d_\infty)$ is not complete:
        
        Take $j \mapsto x_k^{(j)}$ and set $x_k^{(j)} := \begin{cases}
            \frac{1}{k} & 1 \leq j \leq k \\
            0 & \text{else}
        \end{cases}$. For any $m > j$, we get $\abs{x_k^{(j)} - x_k^{(m)}} \leq \frac{1}{j}$, so $d_\infty(x^{(j)},x^{(m)}) \leq \frac{1}{j}$ for $m > j$, so the sequence is Cauchy but has no limit. 
    \end{enumerate}
\end{beispiel}

\begin{prop}
    Let $(X,d)$ be a complete metric space and $A \subseteq X$ closed. Let $d_A := d|_{A \times A}$. Then, $(A, d_A)$ is complete.
\end{prop}
\begin{proof}
    If $x: \IN \to A$ is a Cauchy sequence in $X$, by completeness, $x_k \to x^*$ in $X$. As $A$ is (sequentially) closed, $x^* \in A$.
\end{proof}

\begin{defi}{Isometric immersion, isometry}{isometry}\index{isometry}\index{isometric immersion}
    Let $(X, d_X)$ and $(Y, d_Y)$ be metric spaces. A map $f: X \to Y$ is called an \emph{isometric immersion} if 
    \begin{equation}
        d_Y(f(x),f(y)) = d_X(x,y)
    \end{equation}
    for all $x, y \in X$. The map $f$ is called an \emph{isometry} if it is in addition bijective
\end{defi}

\begin{bemerkung}
    An isometric immersion is automatically bijective. Thus, if $f$ is an isometric immersion, than $f$ is an isometry from $X$ to $f(X)$. The spaces $X$ and $f(X)$ (equipped with the restriction of the metric $d_Y$) are indistinguishable as metric spaces. In particular, a sequence $f: \IN \to \IR$ is convergent in $X$ if and only if $f \circ x$ is convergent in $f(X)$, and $X$ is a Cauchy sequence in $X$ if and only if $f \circ x$ is a Cauchy sequence in $f(X)$.
\end{bemerkung}

\begin{thm}{Completion}{completion}
    Let $(X,d)$ be a metric space. Then, there exists a complete, metric space $(\tilde{X}, \tilde{d})$ and an isometric immersion $j: X \to \tilde{X}$ such that $j(X)$ is dense in $\tilde{X}$.
\end{thm}
The space $\tilde{X}$ is called the \emph{completion of $X$}\index{space!completion}\index{completion} (and is unique up to isometries).

\begin{proof}
    We will see a short proof later. Here we just sketch the standard proof which is based on considering the space of Cauchy sequences in $X$ modulo converging sequences. This proof is modelled on the construction of the real numbers $\IR$ from the rational numbers $\IQ$ (see e.g. Alt's book for details).

    Let
    \begin{equation}
        \hat{X} := \set{x: \IN \to X \mid \text{$x$ Cauchy sequence}}.
    \end{equation}
    Define 
    \begin{equation}
        x \sim y \quad :\iff \quad \lim_{k \to \infty} d(x_k, y_k) = 0.
    \end{equation}
    It is easy to see that $\sim$ is an equivalence relation. The equivalence class of $x \in \hat{X}$ is defined by 
    \begin{equation}
        [x] := \set{z \in \hat{X} \mid z \sim x},
    \end{equation}
    and the space of equivalences is defined by $\hat{X} := \nicefrac{X}{\sim}$. On $\hat{X}$, we define 
    \begin{equation}
        \tilde{d}([x],[y]) := \lim_{j \to \infty} d(x_j, y_j).
    \end{equation}
    It follows from the definition of the equivalence classes and the triangle inequality that $\tilde{d}$ is well-defined, i.e., the right hand side only depends on the equivalence classes of the Cauchy sequences $x$ and $y$. Since $x$ and $y$ are Cauchy sequences, it is also easy to see that $j \mapsto d(x_j, y_j)$ is a Cauchy sequence in $\IR$ and that therefore $\lim_{j \to \infty} d(x_j, y_j)$ exists. It is easy to see that $\tilde{d}$ is a metric on $\tilde{X}$. Then, one shows that
    \begin{itemize}
        \item $(\tilde{X},\tilde{d})$ is complete,
        \item the map $j: X \to \tilde{X}$ which maps $a$ to the (equivalence class of the) constant sequence with value $a$ is an isometry
        \item $j(X)$ is dense in $(\tilde{X},\tilde{d})$
    \end{itemize}
\end{proof}

\begin{lem}{Hausdorff metric space is complete}{hausdorffSpaceComplete}
    Let $(X,d)$ be a complete, metric space. Then, the space $(\cC, d_H)$ of closed, bounded, non-empty set with the Hausdorff metric (introduced in Lemma~\ref{lem:hausdorffDistanceMetric}) is complete.
\end{lem}
\begin{proof}
    Exercise. \emph{Hint:} First show that a decreasing Cauchy sequence of sets $B_j$ converges to $B^* := \bigcap_{j \in \IN} B_j$ (and in particular, $B^* \neq \emptyset$). By passage to a subsequence, you can assume without loss of generality that $d_H(B_j, B_{j+1}) < 2^{-j}$ (explain why). For a general Cauchy sequence $A_j$ define $B_j := \overline{\bigcup_{k \geq j} A_j}$ and show that $d_H(A_j, B_j) \to 0$.
\end{proof}

One can also introduce a metric on set of (equivalence classes of) bounded and complete metric spaces, the so called \emph{Gromov-Hausdorff metric}\index{metric!Gromov-Hausdorff -}\index{Gromov-Hausdorff metric}. Here we say that $(X, d_X) \sim (Y, d_Y)$ if there exist an isometry $I : X \to Y$. Then, one defines
\begin{equation}
    \begin{split}
        &d_{GH}((X,d_X),(Y,d_Y)) \\
        :=& \inf\set{d_{H,Z}(I(X),J(Y)) \mid \text{$(Z,d_Z)$ metric space, $I: X \to Z, J: Y \to Z$ isom. immers.}}
    \end{split}
\end{equation}
where $d_{H,Z}$ denotes the Hausdorff distance in the metric space $(Z,d_Z)$. The GH metric plays an important role in geometry as well as in application to image processing.

\newLecture{18.10.2024}

\newLecture{23.10.2024}



\end{document}